%-------------------------
% Resume in LateX
% Author : Catherine Gilbert
% License : MIT
%------------------------

\documentclass[11pt]{article} % fontsize 10pt/11pt/12pt

\usepackage[margin=1in, letterpaper]{geometry}
\usepackage[colorlinks=true, urlcolor=blue]{hyperref}
\setcounter{secnumdepth}{0} % remove section numbering
\usepackage{titlesec}
% remove vertical spacing above and below
\titlespacing{\subsection}{0pt}{*0}{*0}
\titlespacing{\subsubsection}{0pt}{*0}{*0}
\titleformat{\section}{\large\bfseries\uppercase}{}{}{}[\titlerule]
\titleformat*{\subsubsection}{\large\itshape}
\usepackage{enumitem}
\setlist[itemize]{noitemsep,left=0pt .. \parindent}
\pagestyle{empty} % remove page number
\pdfgentounicode=1


\begin{document}
\begin{center}
	\begin{minipage}{0.5\textwidth}
		{\Huge\bfseries
			Catherine Gilbert
		} \\ \medskip
	\end{minipage} \hfill
	\begin{minipage}{0.4\textwidth}
		\raggedleft
		Email~: recrutement@cathgilbert.com \\
		Cellulaire~: +1 (418) 313-1402 \\
		Site Web~: \url{https://olaren.dev} \\
	\end{minipage}
\end{center}


\section{Expérience}
%=================%
\subsection{Desjardins \hfill Lévis/Montréal, Qc}
\subsubsection{Analyste-développeuse \hfill 2021~---~Présent}
\begin{itemize}
    \item Réduction du temps de chargement de page par 7.5x en parallélisant
        les requêtes de données essentielles et en ajoutant du squelette de
        chargement pour le contenu non critique.
    \item Direction des décisions de conception et d'architecture techniques
        à travers plusieurs projets sur la plateforme AccèsD.
    \item Migration de fonctionnalités bancaires cœur d'un monolithe WebLogic
        legacy vers des microservices Spring Boot.
    \item Développement full-stack sur un backend-for-frontend Spring Boot
        et un frontend Angular.
    \item Chiffrement de données confidentielles à travers la plateforme sans
        impact sur la disponibilité des services.
    \item Mise en place de tests de régression visuelle réduisant la charge
        d'assurance qualité de 20~\%.
\end{itemize}

\subsection{N'Ware Technologies \hfill Saint-Georges, Qc}
\subsubsection{Développeuse stagiaire \hfill Été 2020}
\begin{itemize}
	\item Réécriture et modernisation d'une application de gestion d'entrepôt.
\end{itemize}


\section{Compétences et intérêts}
%===========================%
\begin{description}[itemsep=0pt]
	\item[Programmation] TypeScript\,; Java\,; C\#\,; Rust\,; C++\,; C\,;
	\item[Logiciels] GNU+Linux\,; Neovim\,; Git\,;
	\item[Langues] Français et anglais. Fluente à l'oral et à l'écrit.
	\item[Intérêts] logiciels libres\,; sécurité et vie privée\,;
\end{description}


\section{Projets}
%===========================%
\begin{itemize}
	\item Homelab~: Infrastructure auto-hébergée à usage personnel et familial
		\begin{itemize}
			\item Héberge plusieurs services (Syncthing, Immich,
			Radicale, Forgejo, Jellyfin, Navidrome, ...)
			\item Sauvegardes 3-2-1~: snapshots ZFS locales,
			réplication on-site et à distance. Tests de restauration périodiques.
			\item Sécurité avec SSO sur chaque service (kanidm),
			CA personnalisé pour certificats TLS (step-ca),
			segmentation VLAN et Wireguard.
			\item Stack de surveillance complet avec Prometheus,
			Loki, Alertmanager et Grafana.
		\end{itemize}
	\item diffusion~: Engin de rendu graphique en C99 sans librairie tierce.
	\item maptoposter~: Génération d'affiches de carte de villes qui mettent
        en évidence les réseaux de transports en commun.
	\item ffxiv-status-contrast \& ffxiv-easy-translate~: Plugins pour le jeu
        vidéo Final Fantasy XIV.
	\item hudkit-wayland~: Overlay pour créer des HUD wayland avec les
        technologies du web.
\end{itemize}

\pagebreak{}


\section{Contributions Open Source}
%===========================%
\begin{itemize}
	\item jellyfin-androidtv~: Client Android TV pour le logiciel de serveur
        multimédia Jellyfin.
	\item PDSls~: Explorateur de données sur le protocole AT.
	\item Lumina~: Framework pour interagir avec les données du jeu FFXIV.
	\item bevy~: Engin de jeu ECS écrit en Rust.
\end{itemize}


\section{Éducation}
%=================%
\subsection{CÉGEP Beauce-Appalaches \hfill Saint-Georges, Qc}
\subsubsection{Technique de l'informatique \hfill 2018~---~2021}
\begin{itemize}
	\item Cote R~: 32.8
\end{itemize}

\end{document}
